Retrieval Augmented Generation (RAG) is a key component for generating accurate and hallucination free answers using Large Language Models (LLMs). LLMs are improving at handling long context, but still suffer from ``lost in the middle'' problem. Thus, precise and accurate retrieval is important. Current retrievers chunk long context into length-based manageable chunks -- in the process throwing away rich and informative semantic global structure in the corpus. We introduce a novel retrieval system {\sys} that empowers an LLM to exploit global structure in a corpus such as a Table of Contents ({\toc}) to efficiently store and retrieve information from its model parameters. Our thorough and careful ablation studies with a finetuned Differentiable Search Index (DSI) system show that {\toc} helps build a low hallucination (less than 0.05\%) generative Information Retrieval (IR) system and can generalize to examples where very few training samples are available. To further research in this novel direction of {\toc} based retrieval we release {\ben} -- a diverse benchmark created from 18 books across 6 diverse domains to further research in this novel direction. {\sys} achieves a high Recall@1 score of 82.6\% on {\ben} as compared to DSI (76.9\%), where the difference is found to be statistically significant. {\sys} easily beats other strong baselines such as BM25 (59.5\%), DPR (68.7\%) and out-of-the-box Mistral (13.8\%). The benchmark data and code used for training {\sys} is available at ~\url{https://anonymous.4open.science/r/s_331/README.md}.

% which exploits the rich  in a corpus. {\sys} mimics how humans leverage the structure of a book by using {\toc} to search for information. {\sys} demonstrates that LLMs can be finetuned to build {\toc} based retrievers where the most granular {\toc} entry (leaf node) serves as the retrieval unit. We further release a high quality and novel comprehensive benchmark {\ben} featuring 18 books across 6 diverse domains (with train, test and dev splits) and show that {\sys} achieves Recall@1 score of 82.6\% as compared to the next best system (DSI) (76.9\%) and other baselines such as BM25 (59.5\%), DPR (68.7\%) and out-of-the-box {\basellm} (13.8\%) on average. {\sys} is capable of incorporating corpus knowledge while learning to select the most relevant {\toc} leaf node, resulting in almost zero hallucinations (less than 0.05\%) -- a key issue with LLM-based retrievers. The benchmark data and code used for training {\sys} is available at ~\url{https://anonymous.4open.science/r/s_331/README.md}.

% Modern day retrievers even with latest advancements have unsolved challenges in handling long context, a major one being forced to split long context into chunks which are independently indexed and used as sub-optimal units of retrieval. Following the intuition from previous works showing how a global structure can help in tasks such as knowledge ingestion, we posit that adopting a structured global view over the whole search corpus such as Table of Contents ({\toc}) can act as a very good prior for retrievers too. We pose a new paradigm of retrieval: {\toc} based retrieval and demonstrate through experiments that {\toc} indeed helps to navigate the search corpus more efficiently for precise information retrievals. Given the remarkable success of modern day LLMs in understanding human instructions across a diverse set of tasks, we leverage this inherent capability of LLMs to design a novel retriever {\sys}, as the first system on this Table of Contents ({\toc}) based retrieval paradigm. {\sys} mimics how humans leverage the structure of a book (or a long document) by using {\toc} to search for information. {\sys} demonstrates that LLMs can be finetuned to build {\toc} based retrievers where the most granular {\toc} entry (leaf node) serves as the retrieval unit. We also introduce a novel comprehensive benchmark {\ben} featuring 18 books across 6 diverse domains (with train, test and dev splits) and show that {\sys} achieves Recall@1 score of 82.6\% as compared to the next best system (DSI) (76.9\%) and other baselines such as BM25 (59.5\%), DPR (68.7\%) and out-of-the-box {\basellm} (13.8\%) on average. {\sys} is capable of incorporating corpus knowledge while learning to select the most relevant {\toc} leaf node, resulting in almost zero hallucinations (less than 0.05\%) -- a key issue with LLM-based retrievers. The benchmark data and code used for training {\sys} is available at ~\url{https://anonymous.4open.science/r/s_331/README.md}.